\chapter{Il modulo \code{autoscaler}}
\label{cap:11-autoscaler}

Il modulo \code{autoscaler} decide \emph{quante repliche} deve avere una
funzione gestita. \`E la controparte, e la controfigura, del modulo trattato nel
capitolo successivo: entrambi regolano la capacit\`a offerta a una funzione, ma
agiscono su leve diverse e --- questo \`e il punto architetturale di \nf{} ---
sono sostituibili indipendentemente.

Sono circa 850 righe in undici classi.

\section{Il posto dell'autoscaler nel sistema}

Il campo \code{ScalingStrategy} nella specifica di una funzione ammette tre
valori, e definisce chi possiede la decisione:

\begin{center}
\small
\begin{adjustbox}{max width=\textwidth}
\begin{tabular}{@{}lp{9.4cm}@{}}
\toprule
\textbf{Strategia} & \textbf{Chi decide} \\
\midrule
\code{INTERNAL} & Questo modulo. Il control plane campiona le proprie metriche e
comanda il numero di repliche al provider di deployment. \\
\code{HPA} & Kubernetes. Il control plane crea le risorse e si astiene; la
decisione \`e dell'Horizontal Pod Autoscaler~\cite{k8shpa}. \\
\code{NONE} & Nessuno. Il numero di repliche resta quello impostato
manualmente. \\
\bottomrule
\end{tabular}
\end{adjustbox}
\end{center}

Il default per una funzione in modalit\`a \code{DEPLOYMENT} \`e \code{INTERNAL}
(\cref{cap:06-nucleo}). La strategia \code{NONE} non \`e una lacuna ma uno
strumento sperimentale: consente di fissare il numero di repliche mentre si
studia il comportamento del controllore di concorrenza, isolando le due
dinamiche che altrimenti interagirebbero.

\section{La decisione}

\subsection{La formula}

\code{ScalingDecisionCalculator} implementa la stessa formula dell'HPA di
Kubernetes: il rapporto fra valore osservato e valore obiettivo, moltiplicato
per il numero corrente di repliche.

\[
  R_{\text{racc}} = \left\lceil \max_{m \in M}
    \frac{v_m}{t_m} \cdot R_{\text{corr}} \right\rceil,
  \qquad
  R_{\text{des}} = \operatorname{clamp}(R_{\text{racc}}, R_{\min}, R_{\max})
\]

dove $M$ \`e l'insieme delle metriche configurate, $v_m$ il valore osservato e
$t_m$ l'obiettivo. Con pi\`u metriche vince la pi\`u esigente: il massimo dei
rapporti, non la media. \`E la scelta corretta per un insieme di vincoli che
devono valere tutti. Un obiettivo assente o non numerico ricade su 50.

Un dettaglio del codice attuale cambia il significato di $R_{\text{corr}}$: il
moltiplicatore \`e il numero di repliche \emph{pronte} (\emph{serving}), non quello
richiesto al deployment. Il rapporto misura quanto il carico eccede la capacit\`a
che sta effettivamente servendo; moltiplicarlo per repliche non ancora pronte
sovrastimerebbe la capacit\`a e frenerebbe la risposta a un picco.

\subsection{Una separazione che sembra ridondante}

\begin{lstlisting}[language=Java]
// Kept apart on purpose: the clamp used to overwrite the recommendation in
// place, so a decision pinned at maxReplicas was indistinguishable from a
// metric that had stopped rising.
int recommendedReplicas = (int) Math.ceil(maxRatio * normalizedCurrentReplicas);
int desiredReplicas =
        Math.clamp(recommendedReplicas, scaling.minReplicas(), scaling.maxReplicas());
\end{lstlisting}

Il valore raccomandato e quello desiderato sono conservati entrambi nel record
di decisione, anche se solo il secondo viene applicato. Il commento del record
lo dice esplicitamente: \`e una distinzione che l'HPA di Kubernetes non espone,
perch\'e pubblica soltanto il valore limitato. La ragione \`e
diagnostica: una funzione ferma a dieci repliche pu\`o esserlo perch\'e la
domanda si \`e stabilizzata oppure perch\'e la domanda continua a crescere ma il
tetto \`e stato raggiunto. Sono due situazioni che richiedono azioni opposte
dall'operatore, e sovrascrivendo la raccomandazione con il valore limitato
diventano indistinguibili.

\`E un esempio del principio ricorrente in questo progetto: \emph{conservare la
distinzione fra ci\`o che il sistema voleva e ci\`o che ha potuto}. Lo stesso
principio governa i due limiti di concorrenza del \cref{cap:09-async-queue} e i
due campi \code{requested}/\code{effective} della modalit\`a di esecuzione nel
\cref{cap:07-modalita-esecuzione}.

\section{Le tre metriche}

\subsection{\code{queue\_depth} e \code{in\_flight}}

Entrambe sono lette da una \code{WorkloadMetricsSource}, interfaccia del progetto
\code{workload-metrics} con quattro letture (profondit\`a, in volo, concorrenza
effettiva, backlog smaltibile). L'implementazione non appartiene a un modulo di
coda: \`e \code{EngineWorkloadMetricsSource}, nel control plane, sopra il motore di
scheduling composto (\cref{cap:09-async-queue}). La profondit\`a \`e il numero di
prenotazioni che la funzione occupa nel motore, mantenuto incrementalmente e
letto in tempo costante, in modo che uno scrape non scandisca mai il backlog; il
commento nel sorgente dichiara anche che differisce di al pi\`u $+1$, globalmente,
dalla definizione ``accodato e non ancora reclamato'' dei vecchi scheduler, e che
n\'e il rapporto dell'autoscaler n\'e la soglia del governatore possono essere spostati
da una lettura di $\pm 1$. Il valore in volo \`e quello dei lease di capacit\`a
della funzione. Sono metriche di \emph{stato}: descrivono la situazione
istantanea del sistema.

\subsection{\code{rps}: un caso a parte}

La terza metrica non viene dalla sorgente ma dalle osservazioni sulle
invocazioni (\code{InvocationObservations}), un'interfaccia dell'SPI implementata
da \code{Metrics}, che restituisce uno snapshot con il totale cumulativo dei
dispatch \emph{e la generazione della funzione} a cui appartiene. La lettura la
trasforma in derivata:

\begin{lstlisting}[language=Java]
var sample = observations.snapshot(functionName);
if (sample.generation() == null) { lastDispatchSamples.remove(functionName); return 0; }
CounterSample current = new CounterSample(sample.generation(), sample.dispatched(), now);
CounterSample previous = lastDispatchSamples.put(functionName, current);
if (previous == null || !previous.generation().equals(current.generation())) { return 0.0; }
double deltaCount = current.count() - previous.count();
long deltaMs = current.epochMs() - previous.epochMs();
if (deltaCount <= 0.0 || deltaMs <= 0L) { return 0.0; }
return deltaCount / (deltaMs / 1000.0);
\end{lstlisting}

Il totale \`e cumulativo, e usarne il valore grezzo come carico sarebbe un errore
--- un sistema attivo da un'ora avrebbe un ``carico'' enorme e crescente. La
lettura ne calcola quindi la differenza fra campioni successivi divisa per il
tempo trascorso. Restituisce zero al primo campione, perch\'e non esiste un
precedente, e \emph{al cambio di generazione}: una funzione cancellata e
ri-registrata riparte con contatori nuovi, e sottrarre il totale della vecchia
generazione da quello della nuova produrrebbe una derivata priva di senso
(negativa, quindi azzerata, o spuria). La versione precedente leggeva un
\code{Counter} Micrometer per nome e non aveva questo confine.

La documentazione operativa del progetto lo enuncia esplicitamente: \emph{``The
raw cumulative counter value is not used directly as load.''} \`E un'avvertenza
che vale la pena avere scritta, perch\'e l'errore \`e comune e le sue
conseguenze --- scale-up monotono e irreversibile --- sono spettacolari.

\subsection{Che cosa succede senza un modulo di coda}

\begin{lstlisting}[language=Java]
@ConditionalOnBean({WorkloadMetricsSource.class, FunctionCatalogView.class})
public class AutoscalerConfiguration { ... }
\end{lstlisting}

L'autoscaler \`e condizionato all'esistenza di una \code{WorkloadMetricsSource}, e
questa esiste solo se il motore di scheduling esiste, cio\`e se almeno un modulo di
coda fornisce una strategia. Senza, il modulo \emph{non viene creato}, e non c'\`e pi\`u
la degradazione parziale (``\code{rps} funziona, il resto no'') che le versioni
precedenti segnalavano con un avviso emesso una volta. Il commento nel sorgente
dichiara la lezione di un incidente sullo stesso meccanismo: la condizione
elencava anche il \code{MeterRegistry}, valutata prima che la sua definizione
esistesse, falliva e portava con s\'e l'intero modulo \emph{in silenzio} --- nessun
bean, nessun \code{InternalScaler}, nessuna riga di log --- e una campagna ha
proseguito con 340 richieste al secondo contro una soglia di 100 restando su una
replica. Il registro \`e oggi iniettato nei bean e non figura nella condizione,
cos\`i che una sua assenza sarebbe rumorosa. Resta per\`o un caso silenzioso, che
il commento di \code{EngineWorkloadMetricsSource} nomina: senza la sorgente,
autoscaler e governatore di concorrenza sono entrambi disattivati senza errore.
Chi vuole l'autoscaler deve quindi selezionare almeno un modulo di coda.

\section{Isteresi: i cooldown}

Un controllore reattivo che agisce a ogni campione oscilla. \`E il fenomeno che
in teoria del controllo si chiama \emph{hunting}: la correzione modifica la
grandezza osservata, che genera la correzione opposta.

Il modulo lo previene con due cooldown asimmetrici:

\begin{lstlisting}[language=Java]
private static final long SCALE_UP_COOLDOWN_MS = 30_000;
private static final long SCALE_DOWN_COOLDOWN_MS = 60_000;
\end{lstlisting}

L'asimmetria \`e il punto. Scalare in su \`e un'operazione a basso rischio: si
paga capacit\`a inutilizzata per qualche minuto. Scalare in gi\`u \`e ad alto
rischio: se la domanda ritorna, le richieste pagano un cold start. Quindi il
sistema sale rapidamente e scende con riluttanza --- lo stesso principio che
governa le finestre di stabilizzazione dell'HPA di Kubernetes, dove il default di
scale-down \`e di cinque minuti.

I valori 30\,s e 60\,s sono costanti nel codice, non configurabili
(\code{ScalingCooldownTracker}); sono configurabili solo l'intervallo di
polling e i limiti di default di repliche (\code{nanofaas.scaling.*}). \`E una
scelta di semplicit\`a che il progetto potrebbe voler rivedere: sono i due
parametri che pi\`u probabilmente un esperimento vorrebbe variare.

\begin{damisurare}
L'effetto dei due cooldown sulla stabilit\`a del numero di repliche e sul costo
in capacit\`a inutilizzata non \`e stato caratterizzato. Le campagne del
\cref{cap:23-prestazioni-base} riportano il numero di repliche di picco per
release ma non la traiettoria temporale.
\todomis{traiettoria del numero di repliche sotto carico a gradino, al variare
dei cooldown}
\end{damisurare}

\section{Rollout in stallo e protezione dal risveglio}

Il ciclo confronta la raccomandazione non con le repliche pronte ma con quelle
\emph{gi\`a richieste} al deployment, e lo fa leggendo richieste e pronte dallo
\emph{stesso} snapshot, perch\'e una decisione non sia presa su un target che nel
frattempo si \`e mosso. Ne segue un problema che il codice risolve con un
componente proprio. Se la raccomandazione scende sotto il richiesto mentre il
rollout sta ancora raggiungendo il richiesto, tornare indietro subito
annullerebbe uno scale-up che sta funzionando (il commento rimanda a un test di
regressione, \code{InternalScalerDesiredVsReadyRegressionTest}); ma non pu\`o
essere una protezione senza fine, perch\'e se alcune repliche non diventano mai
pronte un vero scale-down resterebbe bloccato per sempre su un target fantasma.
\code{ScalingProgressTracker} decide quando ``ancora in salita'' diventa
``bloccato'': il rollout \`e in stallo se il richiesto non \`e raggiunto e il numero di
pronte non \`e avanzato per una finestra intera di 60\,s.

Lo scale-down segue quindi due strade: il \emph{segnale di scale-down} esplicito
(le repliche pronte superano gi\`a ci\`o che il carico richiede, e non attende il
completamento del rollout), oppure lo stallo. In entrambi i casi passa per
\code{scaleDownIfUnprotected} del coordinatore di risveglio, che rifiuta se la
generazione ha un lease di risveglio attivo (\cref{sec:06-riprova}): non si toglie
capacit\`a a una funzione che una richiesta sta aspettando. La lettura delle
repliche \`e un'osservazione non bloccante con tre esiti (fresca, vecchia,
assente); senza alcuna lettura il ciclo salta la funzione, perch\'e trattarla come
zero repliche scalerebbe un deployment sano sulla base di una GET fallita.

\section{Inferenza del cold start}

\code{ColdStartTracker} risolve un problema di osservabilit\`a: un runtime
vecchio, o scritto in un linguaggio il cui SDK non implementa la convenzione,
non emette l'header \code{X-Cold-Start}, e il control plane non saprebbe che
un'invocazione ha pagato un'inizializzazione.

Il tracker ricorda gli eventi di scale-up recenti e ne inferisce il cold start:

\begin{lstlisting}[language=Java]
/**
 * Returns true if the function recently scaled up from 0 replicas
 * (within the expiry window), indicating a likely cold start.
 */
public boolean isPotentialColdStart(String functionName) {
    ScaleUpEvent event = recentScaleUps.get(functionName);
    if (event == null) { return false; }
    long ageMs = Instant.now().toEpochMilli() - event.timestamp().toEpochMilli();
    if (ageMs > EXPIRY_MS) { recentScaleUps.remove(functionName); return false; }
    return event.fromReplicas() == 0;
}
\end{lstlisting}

\`E un'euristica, e il codice non lo nasconde --- il metodo si chiama
\code{isPotentialColdStart}, non \code{isColdStart}. La finestra di 30 secondi \`e
generosa, e il criterio \code{fromReplicas() == 0} \`e conservativo: uno scale-up
da una a due repliche non \`e considerato, anche se la seconda replica sar\`a
fredda.

L'euristica pu\`o sbagliare in entrambe le direzioni: attribuisce cold start a
richieste servite da una replica gi\`a calda entro trenta secondi da un
risveglio, e non lo attribuisce a richieste servite da una replica nuova
aggiunta a un insieme non vuoto. \`E accettabile per il ruolo che ha --- un
ripiego per runtime che non collaborano --- e diventa problematica se qualcuno
ne usasse i risultati come dato sperimentale.

\begin{damisurare}
Il tasso di errore dell'inferenza rispetto all'header \code{X-Cold-Start}
riportato dai runtime che lo implementano \`e misurabile confrontando le due
sorgenti sulla stessa esecuzione, e non \`e stato misurato.
\todomis{accordo fra cold start inferito e cold start dichiarato}
\end{damisurare}

\section{Il ciclo di scaling}

\code{InternalScaler} \`e un \code{SmartLifecycle} con un esecutore schedulato a
periodo fisso (default 5\,s, \code{nanofaas.scaling.poll-interval-ms}). A ogni tick
scandisce le funzioni registrate, e per ciascuna con strategia \code{INTERNAL},
deployment gestito e generazione corrente calcola la decisione, verifica il
cooldown pertinente, e comanda il nuovo numero di repliche al coordinatore di
deployment. Un'eccezione su una funzione non ferma le altre.

Il modulo si disattiva completamente in assenza di un coordinatore di deployment:

\begin{lstlisting}[language=Java]
if (deploymentCoordinator == null) {
    log.info("InternalScaler disabled: no ManagedReplicaControl available");
    return;
}
\end{lstlisting}

\`E la degradazione corretta: senza un backend di deployment non esistono
repliche da scalare, e un ciclo che gira senza poter agire sarebbe puro consumo.
Anche qui la condizione \`e dichiarata nel log all'avvio anzich\'e lasciata
dedurre.

Il periodo di 5\,s \`e sostanzialmente pi\`u lungo dei tempi dell'ammissione, e la
differenza \`e istruttiva: l'ammissione deve reagire in millisecondi perch\'e
governa singole richieste, lo scaling reagisce in secondi perch\'e la sua azione
--- creare un pod --- richiede secondi. Un autoscaler che campionasse a pochi
millisecondi produrrebbe decisioni su uno stato che non ha ancora recepito
l'azione precedente.

\section{Metriche di decisione}

\code{ScalingDecisionMetrics} esporta la decisione stessa, in quattro gauge per
funzione: repliche raccomandate, desiderate, un indicatore di limitazione
(\code{function\_scaling\_limited}) e il rapporto massimo moltiplicato per 1\,000
(\code{function\_scaling\_ratio\_milli}, intero perch\'e un rapporto di 3,65 non
sia riportato come 3 o 4). \code{TargetLoadMetrics} esporta invece gli obiettivi
configurati (\code{gateway\_service\_target\_load}) in una forma compatibile con
OpenFaaS. Rende
possibile ricostruire \emph{a posteriori} perch\'e il sistema abbia scelto un
certo numero di repliche, senza dover correlare i log con le metriche di stato.

La documentazione operativa fornisce la lettura combinata: \emph{``Internal
autoscaling for \code{rps} uses the delta between successive
\code{function\_dispatch\_total} samples divided by elapsed sample time''}, ed \`e
la stessa quantit\`a che l'operatore ricalcolerebbe da PromQL. Esportare la
grandezza su cui il sistema ha effettivamente deciso, anzich\'e lasciarla
ricostruire, evita la classe di errori in cui l'analisi post-mortem calcola
qualcosa di leggermente diverso da ci\`o che il controllore aveva visto.

\section{Rapporto con il controllo della concorrenza}

I due moduli agiscono su leve diverse:

\begin{center}
\small
\begin{adjustbox}{max width=\textwidth}
\begin{tabular}{@{}lll@{}}
\toprule
 & \textbf{\code{autoscaler}} & \textbf{\code{concurrency-control}} \\
\midrule
Leva & numero di repliche & limite di concorrenza per funzione \\
Latenza dell'azione & secondi (creazione pod) & immediata (variabile in memoria) \\
Periodo di decisione & 5\,s & \multicolumn{1}{l}{\`e nel \cref{cap:12-concurrency-control}} \\
Costo dell'errore & capacit\`a sprecata o cold start & attesa o sotto-utilizzazione \\
\bottomrule
\end{tabular}
\end{adjustbox}
\end{center}

Il modo \code{STATIC\_PER\_POD} del controllore di concorrenza \`e l'unico punto
in cui i due si toccano: calcola il limite come \emph{repliche pronte $\times$
obiettivo per replica}, e quindi il suo risultato si muove quando l'autoscaler
agisce. \`E un accoppiamento voluto e unidirezionale --- il controllore di
concorrenza legge lo stato delle repliche, l'autoscaler non legge i limiti di
concorrenza.

Che i due possano essere selezionati indipendentemente \`e il fatto
architettonico che rende possibili gli esperimenti del
\cref{cap:25-valutazione-concorrenza}: fissando la strategia di scaling a
\code{NONE} si ottiene un sistema in cui l'unica variabile \`e il limite di
concorrenza, e le due dinamiche non si confondono.
